\documentclass{article}
\usepackage[UTF8]{ctex}
\usepackage{amsmath,amssymb,amsfonts}
\usepackage{enumitem}
\usepackage{geometry}
\geometry{a4paper, margin=1in}
% === 错题本排版优化（不参与题目标识） ===
\setlength{\parindent}{0pt}
\setlength{\parskip}{0.35em plus 0.10em minus 0.08em}
\setlength{\emergencystretch}{3em}
\setlist{itemsep=0.25em, topsep=0.40em, parsep=0pt}
\allowdisplaybreaks[2]
% === 错题本排版优化结束 ===
\begin{document}
\section*{浙江大学2026年数学分析真题}
\section*{103．浙江大学2026年数学分析试卷}

\subsection{题目1}
1．设 \(x_{n}=\frac{1}{n+1}+\frac{1}{n+2}+\cdots+\frac{1}{n+n}, n \in \mathbb{Z}^{+}\)，求极限 \(\lim _{n \rightarrow \infty} x_{n}\) 。

\subsection{题目2}
2．设 \(f(x)\) 在区间 \((0,2)\) 上可导，且 \(f(1)=f^{\prime}(1)=0, f^{\prime \prime}(1)=2\) ，求极限 \(\lim _{x \rightarrow 1} \frac{x f^{\prime}(x)}{x f^{\prime}(x)-f(x)}\) ．

\subsection{题目3}
3．设 \(x_{n}=\int_{0}^{\frac{\pi}{2}} \frac{\sin ^{2}(n x)}{\sin x} \mathrm{~d} x, n \in \mathbb{Z}^{+}\)，求极限 \(\lim _{n \rightarrow \infty} \frac{x_{n}}{\ln n}\) 。

\subsection{题目4}
4．设 \(S\) 是 \(x^{2}+y^{2}=z^{2}\) 介于 \(z=0\) 与 \(z=1\) 之间的部分，并取外侧为正向，计算第二型曲面积分

\[
\iint_{S}(x-2 y+z) \mathrm{d} y \mathrm{~d} z+(y-2 z+x) \mathrm{d} z \mathrm{~d} x+(z-2 x+y) \mathrm{d} x \mathrm{~d} y .
\]

二．（15 分）设 \(f(x)=\left\{\begin{array}{ll}x\left(1-x^{2}\right), & x \in \mathbb{Q}, \\ x\left(1+x^{2}\right), & x \notin \mathbb{Q} .\end{array}\right.\) 请讨论 \(f(x)\) 在 \(\mathbb{R}\) 上的连续性与可导性（需说明理由）．

三．（10 分）证明：设 \(f(x)\) 在区间 \([0,1]\) 上只有第一类间断点，则 \(f(x)\) 是 \([0,1]\) 上的有界函数．

四．（15 分）设 \(f(x)\) 在区间［0，1］上二阶可导，且 \(f(0)=f(1)=0\) ，若 \(\min _{x \in[0,1]} f(x)=-1\) ，证明：存在\(\xi \in(0,1)\) ，使得 \(f^{\prime \prime}(\xi) \geq 8\) ．

五．（15 分）证明：函数 \(y=\ln x\) 在 \([a,+\infty)\) 上一致连续，其中 \(a>0\) ，但在 \((0,+\infty)\) 上不一致连续．

六．（15 分）设 \(a_{1}>0, a_{2}>0\) ，定义 \(a_{n+2}=\sqrt{a_{n+1}}+\sqrt{a_{n}}, n=1,2, \cdots\) ，证明：数列 \(\left\{a_{n}\right\}\) 收敛，并求\(\lim _{n \rightarrow \infty} a_{n}\) ．

七．（15 分）设 \(\lambda \geq 1\) ，定义 \(I(\lambda)=\int_{0}^{1} \cos \left(\lambda x^{2}\right) \mathrm{d} x\) ，证明：存在与 \(\lambda\) 无关的常数 \(c>0\) ，使得\(|I(\lambda)| \leq \frac{c}{\sqrt{\lambda}}\) ．

八．（15 分）设 \(f(x), g(x)\) 是区间 \((0,1)\) 上的有界函数，定义

\[
\varphi(x)=\sup _{t \in(0,1)}(f(t)+\sqrt{x} g(t)), x \in(0,1) .
\]

证明：存在常数 \(C \geq 0\) ，使得对任意 \(x, y \in(0,1)\) ，有 \(|\varphi(x)-\varphi(y)| \leq C \sqrt{|x-y|}\) ．

九．（10 分）证明：函数 \(f(x)=\sum_{n=2}^{\infty} \frac{\sin n x}{n \ln n}\) 在 \(x=0\) 处连续．
\end{document}
